\section{总体设计}

\subsection{系统总体架构}

系统采用前后端分离和容器化部署方式，在线服务链路如下：

\begin{figure}[H]
\centering
\begin{tikzpicture}[
  node distance=0.9cm,
  archbox/.style={
    rectangle,
    rounded corners=2pt,
    draw=black,
    align=center,
    minimum width=4.5cm,
    minimum height=0.85cm,
    font=\song\xiaowuhao
  },
  arrow/.style={-{Stealth[length=2.2mm]}, line width=0.5pt}
]
\node[archbox] (browser) {用户浏览器};
\node[archbox, below=of browser] (frontend) {Streamlit 前端};
\node[archbox, below=of frontend] (backend) {FastAPI 推荐服务};
\node[archbox, below left=0.8cm and 0.5cm of backend] (redis) {Redis 推荐缓存};
\node[archbox, below right=0.8cm and 0.5cm of backend] (postgres) {PostgreSQL 数据库};

\draw[arrow] (browser) -- (frontend);
\draw[arrow] (frontend) -- (backend);
\draw[arrow] (backend) -- (redis);
\draw[arrow] (backend) -- (postgres);
\draw[arrow] (redis) -- (postgres);
\end{tikzpicture}
\caption{系统总体架构}
\end{figure}

其中，Streamlit 负责交互页面展示，FastAPI 负责提供 REST API，PostgreSQL 负责保存电影、用户、
评分和点击事件数据，Redis 用于缓存推荐结果和降低重复计算开销。该架构将展示层、服务层、数据层和
缓存层解耦，便于后续替换推荐算法或扩展部署规模。

\subsection{系统模块划分}

系统主要划分为以下模块：

\begin{center}
\begin{tabular}{ll}
\toprule
\textbf{模块} & \textbf{主要职责} \\
\midrule
backend  & FastAPI 接口、数据库模型、推荐逻辑和系统指标 \\
frontend & Streamlit 页面、电影搜索、推荐结果展示和评估页面 \\
trainer  & 后续 ALS/BPR 等离线模型训练脚本 \\
evaluator & 抽样评估、RecBole 配置和实验结果记录 \\
spark    & Spark 预处理与 Spark MLlib ALS 对照实验入口 \\
docs     & 课程设计报告、实验记录和系统说明文档 \\
\bottomrule
\end{tabular}
\end{center}

\subsection{数据集与数据库设计}

当前系统已经导入 MovieLens 现代 CSV 数据集。数据导入流程首先读取 movies.csv 和 ratings.csv，
将电影信息和评分行为写入 PostgreSQL；如果目录中存在 links.csv 和 tags.csv，系统也可以将外部编号
映射和用户标签导入对应扩展表。导入完成后，推荐服务主要围绕电影表和评分表进行查询与相似度计算。

数据库统计如下：

\begin{center}
\begin{tabular}{lr}
\toprule
\textbf{表} & \textbf{数量} \\
\midrule
users   & 200{,}948 \\
movies  & 87{,}618 \\
ratings & 32{,}921{,}437 \\
\bottomrule
\end{tabular}
\end{center}

数据库核心表包括 users、movies、ratings、user\_events、recommendation\_cache 和
experiment\_results。其中，movies 保存电影标题与类型，ratings 保存用户评分行为，
user\_events 用于记录前端点击事件，recommendation\_cache 用于保存推荐缓存，
experiment\_results 用于记录评估指标。

各类数据在系统中的作用如下：

\begin{center}
\begin{tabular}{ll}
\toprule
\textbf{数据类型} & \textbf{系统作用} \\
\midrule
电影基础信息 & 支持电影搜索、标题展示、类型展示和新增电影演示 \\
用户评分数据 & 支持电影相似度计算、协同过滤推荐和抽样离线评估 \\
用户点击事件 & 记录 Demo 曝光与点击行为，用于计算本地 Demo CTR \\
推荐缓存数据 & 缓存重复推荐请求结果，降低在线计算开销 \\
实验结果数据 & 保存 Precision@K、Recall@K、NDCG@K 和延迟等指标 \\
\bottomrule
\end{tabular}
\end{center}

与 MovieLens 1M 相比，当前数据规模更大，能够更好体现数据库、缓存和服务化部署的必要性。
但由于 MovieLens Latest 数据集会随时间更新，因此正式固定实验结果更适合使用 MovieLens 32M
作为稳定 benchmark。

\newpage
